\emph{Prevalence as a swept parameter.} LSPR23 is a live-fire exercise: $10.06\%$ of flows and
0.48--0.81\% of episodes are
malicious, the optimistic case, so \cref{apptab:prevalence} treats prevalence as a swept parameter rather
than asserting a figure for a typical SOC: the targets run from 0.7 to 2.7 orders of
magnitude below the exercise's episode rate, a sensitivity range, not a claim about any deployment. We thin malicious
episodes --- not flows, which would change group arity and evidence sums and so confound the
feasibility question with the padding cost --- to each target $\pi$, keeping a nested random subset so
a lower target keeps a strict subset of a higher one. The primary mechanism replaces each thinned
episode with benign evidence drawn from the same window, so the hypothesis stays at its own index and
$T$, the margin and every survivor's level are held exactly fixed; the deletion mechanism removes it
instead, which shrinks $T$ by up to 0.80\% and raises the margin by up to
+0.0129. The two agree to within 0.02 in
bootstrap probability, so it is prevalence and not re-indexing that drives the collapse. Feasibility is
untouched while detection collapses: the feasibility-is-not-detection separation, driven by base rate
alone. The $\pi=0$ arm produces no rejection on any stream, but that is 5
effective (overlapping) windows, and zero would be the likely outcome even at the $4.5\%$ early-false
rate quoted from simulation in \cref{app:taxonomy} ($P=0.79$): it is
consistency, not a test of that number.
